\documentclass[12pt,a4paper]{article}

\usepackage[utf8]{inputenc}
\usepackage[T1]{fontenc}
\usepackage[spanish]{babel}
\usepackage{geometry}
\usepackage{xcolor}
\usepackage{listings}
\usepackage{hyperref}
\usepackage{array}
\usepackage{longtable}

\geometry{margin=2.5cm}

\definecolor{codebg}{RGB}{248,248,248}
\definecolor{codeframe}{RGB}{220,220,220}
\definecolor{codekeyword}{RGB}{0,92,175}
\definecolor{codestring}{RGB}{120,70,0}

\lstdefinestyle{pythonstyle}{
    language=Python,
    backgroundcolor=\color{codebg},
    frame=single,
    rulecolor=\color{codeframe},
    basicstyle=\ttfamily\small,
    keywordstyle=\color{codekeyword}\bfseries,
    stringstyle=\color{codestring},
    showstringspaces=false,
    breaklines=true,
    tabsize=4
}

\lstdefinestyle{textstyle}{
    backgroundcolor=\color{codebg},
    frame=single,
    rulecolor=\color{codeframe},
    basicstyle=\ttfamily\small,
    showstringspaces=false,
    breaklines=true
}

\title{Laboratorio 10\\Estilos de Programación RF.1 y RF.2}
\author{Lorenzo \\ Proyecto ChambeaYa}
\date{\today}

\begin{document}

\maketitle

\section{Datos generales}

\textbf{Requisitos bajo responsabilidad:}

\begin{itemize}
    \item RF.1 - Gestión de Registro y Autenticación.
    \item RF.2 - Gestión del Perfil del Joven.
\end{itemize}

\textbf{Recursos de software:}

\begin{itemize}
    \item IDE: VS Code.
    \item Plugin de revisión: SonarLint.
    \item Control de versiones: Git y GitHub.
    \item Gestión del proyecto: Trello con tablero Kanban/Scrum.
\end{itemize}

\section{Contexto del proyecto}

ChambeaYa es una plataforma web para conectar practicantes con empresas y gestionar prácticas preprofesionales. RF.1 permite registrar y autenticar usuarios. RF.2 permite que el joven practicante gestione su perfil profesional.

Según el tablero Trello del proyecto, las funcionalidades de alto nivel incluyen:

\begin{itemize}
    \item RF.1 - Gestión de Registro y Autenticación.
    \item RF.2 - Gestión del Perfil del Joven.
    \item RF.3 - Publicación y Gestión de Convocatorias.
    \item RF.4 - Matching y Postulación.
    \item RF.5 - Desarrollo, Entregables y Evaluación de la Práctica.
    \item RF.6 - Certificación Digital y Reputación.
    \item RF.7 - Notificaciones.
    \item RF.8 - Reportes y Analítica.
\end{itemize}

El prototipo del tablero incluye pantallas relacionadas con registro, perfil profesional y convocatorias. Para este informe se consideran las pantallas de registro y perfil, asociadas directamente con RF.1 y RF.2.

\section{Estructura DDD y arquitectura}

La implementación está organizada por capas:

\begin{lstlisting}[style=textstyle]
domain         -> entidades y reglas de negocio
application    -> casos de uso
presentation   -> controllers
infrastructure -> repositorios
frameworks     -> Flask, rutas, SQLAlchemy y migraciones
\end{lstlisting}

Flujo aplicado en RF.1:

\begin{lstlisting}[style=textstyle]
auth_routes.py
  -> UsuarioController
  -> UsuarioApplicationService
  -> Usuario / AutenticacionDominioServicio
  -> SqlAlchemyUsuarioRepository
  -> UsuarioModel
\end{lstlisting}

Flujo aplicado en RF.2:

\begin{lstlisting}[style=textstyle]
perfil_routes.py
  -> PerfilController
  -> PerfilApplicationService
  -> Practicante
  -> SqlAlchemyPerfilRepository
  -> PracticanteModel
\end{lstlisting}

\section{Estilo 1: RESTful}

\textbf{Práctica:} se exponen recursos mediante endpoints HTTP con verbos alineados a la acción. Se usa \texttt{GET} para consultar, \texttt{POST} para crear o ejecutar acciones y \texttt{PUT} para actualizar.

\textbf{Fragmento RF.1:}

\begin{lstlisting}[style=pythonstyle]
@auth_bp.post("/register")
def register():
    payload = request.get_json(silent=True) or {}
    controller = build_usuario_controller()
\end{lstlisting}

\textbf{Fragmento RF.2:}

\begin{lstlisting}[style=pythonstyle]
@perfil_bp.put("/me")
def update_my_profile():
    usuario_id = get_authenticated_user_id()
    if usuario_id is None:
        return jsonify({"error": "No autenticado"}), 401
\end{lstlisting}

\textbf{Aplicación:}

\begin{itemize}
    \item \texttt{POST /auth/register}: registro de usuario.
    \item \texttt{POST /auth/login}: inicio de sesión.
    \item \texttt{GET /auth/me}: consulta del usuario autenticado.
    \item \texttt{PUT /perfil/me}: creación o actualización de perfil.
    \item \texttt{GET /perfil/me/reputacion}: consulta de reputación.
\end{itemize}

\section{Estilo 2: Error/Exception Handling}

\textbf{Práctica:} las reglas de negocio lanzan excepciones controladas y Flask las transforma en respuestas HTTP claras.

\textbf{Fragmento RF.1:}

\begin{lstlisting}[style=pythonstyle]
try:
    data, status_code = controller.login(payload)
except ValueError as exc:
    return jsonify({"error": str(exc)}), 401
\end{lstlisting}

\textbf{Fragmento RF.2:}

\begin{lstlisting}[style=pythonstyle]
try:
    data, status_code = controller.update_practicante(usuario_id, payload)
except ValueError as exc:
    return jsonify({"error": str(exc)}), 400
\end{lstlisting}

\textbf{Aplicación:}

\begin{itemize}
    \item Credenciales inválidas devuelven \texttt{401}.
    \item Usuario no autenticado devuelve \texttt{401}.
    \item Datos inválidos de perfil devuelven \texttt{400}.
    \item Perfil inexistente devuelve \texttt{404}.
\end{itemize}

\section{Estilo 3: Persistent-Tables}

\textbf{Práctica:} la persistencia se implementa mediante tablas SQLAlchemy y repositorios que encapsulan el acceso a base de datos.

\textbf{Fragmento RF.1:}

\begin{lstlisting}[style=pythonstyle]
class SqlAlchemyUsuarioRepository(IUsuarioRepository):
    def save(self, usuario):
        model = None
        if usuario.id is not None:
            model = db.session.get(UsuarioModel, usuario.id)

        if model is None:
            model = UsuarioModel()
            db.session.add(model)
\end{lstlisting}

\textbf{Fragmento RF.2:}

\begin{lstlisting}[style=pythonstyle]
class SqlAlchemyPerfilRepository(IPerfilRepository):
    def find_practicante_by_user_id(self, usuario_id):
        model = PracticanteModel.query.filter_by(usuario_id=usuario_id).first()
        return self._to_practicante_domain(model)
\end{lstlisting}

\textbf{Aplicación:}

\begin{itemize}
    \item \texttt{UsuarioModel} representa la tabla \texttt{usuarios}.
    \item \texttt{PracticanteModel} representa la tabla \texttt{practicantes}.
    \item Los repositorios aíslan SQLAlchemy de \texttt{domain} y \texttt{application}.
\end{itemize}

\section{Estilo 4: Things}

\textbf{Práctica:} el dominio se modela con objetos que encapsulan estado y comportamiento.

\textbf{Fragmento RF.1:}

\begin{lstlisting}[style=pythonstyle]
class Usuario:
    def login(self):
        if self.estado != self.ESTADO_ACTIVO:
            raise ValueError("La cuenta no está activa")

        return self
\end{lstlisting}

\textbf{Fragmento RF.2:}

\begin{lstlisting}[style=pythonstyle]
class Practicante:
    def verificar_identidad(self):
        if not self.dni and not self.carnet_universitario:
            raise ValueError("Se requiere DNI o carnet universitario")

        self.identidad_verificada = True
        return self
\end{lstlisting}

\textbf{Aplicación:}

\begin{itemize}
    \item \texttt{Usuario} contiene reglas de login, activación y actualización de contraseña.
    \item \texttt{Practicante} contiene reglas de habilidades, formación, identidad y reputación.
\end{itemize}

\section{Estilo 5: Pipeline}

\textbf{Práctica:} las operaciones se implementan como una secuencia clara de pasos: validar, consultar, ejecutar regla de dominio, persistir y responder.

\textbf{Fragmento RF.1:}

\begin{lstlisting}[style=pythonstyle]
def register_user(self, email, password, tipo):
    self._require_repository()

    normalized_email = self._normalize_email(email)
    if self.usuario_repository.find_by_email(normalized_email) is not None:
        raise ValueError("El email ya está registrado")

    password_hash = self.autenticacion_dominio_servicio.generar_password_hash(password)
\end{lstlisting}

\textbf{Fragmento RF.2:}

\begin{lstlisting}[style=pythonstyle]
def register_identity(self, usuario_id, dni=None, carnet_universitario=None):
    self._require_repositories()
    practicante = self._get_or_create_practicante(usuario_id)
    practicante.actualizar_datos(dni=dni, carnet_universitario=carnet_universitario)
    practicante.verificar_identidad()
    practicante.calcular_score()
    return self.perfil_repository.save_practicante(practicante)
\end{lstlisting}

\section{Estilo 6: Cookbook}

\textbf{Práctica:} se usan funciones pequeñas y reutilizables como recetas para construir componentes.

\textbf{Fragmento RF.1:}

\begin{lstlisting}[style=pythonstyle]
def build_usuario_controller():
    repository = SqlAlchemyUsuarioRepository()
    service = UsuarioApplicationService(repository)
    return UsuarioController(service)
\end{lstlisting}

\textbf{Fragmento RF.2:}

\begin{lstlisting}[style=pythonstyle]
def build_perfil_controller():
    perfil_repository = SqlAlchemyPerfilRepository()
    usuario_repository = SqlAlchemyUsuarioRepository()
    service = PerfilApplicationService(perfil_repository, usuario_repository)
    return PerfilController(service)
\end{lstlisting}

\section{Correcciones de calidad}

Sobre RF.1 y RF.2 se aplicaron correcciones compatibles con revisión SonarLint:

\begin{itemize}
    \item Se eliminaron constructores vacíos que solo contenían \texttt{pass}.
    \item Se reemplazaron métodos fuera del alcance de RF.2 con \texttt{NotImplementedError}.
    \item Se evitó exponer contraseñas en respuestas JSON.
    \item Se usó hashing de contraseñas con Werkzeug.
    \item Se usaron tokens aleatorios con expiración para activación y recuperación.
\end{itemize}

\section{Tablero Trello}

El proyecto usa un tablero Trello basado en Kanban/Scrum y User Story Mapping. En las capturas del tablero se observan columnas como:

\begin{itemize}
    \item Product Backlog.
    \item Sprint Backlog.
    \item Prototipo.
    \item Capas del proyecto.
    \item En progreso.
    \item En revisión.
    \item Completado.
\end{itemize}

Historias relacionadas con RF.1 y RF.2:

\begin{itemize}
    \item HF.1.1.1 Registro con email y contraseña.
    \item HF.1.2.1 Login con email/contraseña.
    \item HF.1.2.2 Recuperación de contraseña.
    \item HF.2.1.1 Agregar habilidades técnicas al perfil.
    \item HF.2.1.2 Agregar formación educativa.
    \item HF.2.1.3 Registrar DNI o carnet universitario.
    \item HF.2.1.3 Verificar identidad del practicante.
    \item HF.2.2 Ver y editar perfil.
    \item HF.2.2.1 Consultar score de reputación.
\end{itemize}

\section{Resumen}

\begin{longtable}{|p{4cm}|p{10cm}|}
\hline
\textbf{Estilo} & \textbf{Uso en RF.1/RF.2} \\
\hline
RESTful & Endpoints HTTP para autenticación y perfil. \\
\hline
Error/Exception Handling & Conversión de excepciones de negocio en respuestas HTTP. \\
\hline
Persistent-Tables & Persistencia con SQLAlchemy y tablas \texttt{usuarios}/\texttt{practicantes}. \\
\hline
Things & Entidades \texttt{Usuario} y \texttt{Practicante} con reglas propias. \\
\hline
Pipeline & Flujos secuenciales en servicios de aplicación. \\
\hline
Cookbook & Funciones constructoras reutilizables para controllers. \\
\hline
\end{longtable}

\end{document}
